\section{Experiments and Results}

We evaluate our agents in two distinct regimes: on a set of $1000$ \textit{test tasks} sampled from the same distribution as the training tasks, and on a set of $30$ single-agent and $28$ multi-agent \emph{hand-authored probe tasks}. A rejection sampling procedure guarantees that the procedural test tasks and probe tasks are outside the training set. The probe tasks represent situations that are particularly intuitive to humans, and deliberately cover a wide range of qualitatively different adaptation behaviours. Example probe tasks are depicted in Figures \labelcref{fig:wrong-pair-disappears-explained,fig:pyramid-in-a-haystack-explained,fig:no-touchy-explained} in the Appendix, and a full description of every probe task is available in Appendix \ref{appendix:human_scale_adaptation}. 

The total achievable reward on each task varies, so whenever we present aggregated results on the test or hand-authored task set, we normalise the total per-trial reward for each task against the reward obtained by fine-tuning AdA on the respective task set. We refer to this normalised reward as a \textit{score}. We stipulate that an adaptive agent must have two capabilities: zero-shot generalisation and few-shot adaptation. Zero-shot \textit{generalisation} is assessed by the score in the case of only being given $1$ trial of interaction with a held-out task. Few-shot \textit{adaptation} is assessed by the improvement in score as the agent is given progressively more trials ($k$) of interaction with the task. More precisely, for each $k$ we report the score in the last trial, showing whether or not an agent is able to make use of additional experience on-the-fly to perform better, i.e. measuring adaptation. 

We aggregate scores across a task set using (one or more) percentiles. When presenting individual probe tasks we report unnormalised total last trial rewards per task for agents and for human players where applicable. For full details of our evaluation methodology see Appendix~\ref{app:evaluation}. 

The space of training configurations for AdA is large, comprising model size, auto-curriculum, memory architecture, memory length, number of tasks in the XLand task pool, single vs multi-agent tasks, distillation teacher, and number of training steps. We use a consistent training configuration within each experimental comparison, but different configurations across different experimental comparisons. We therefore caution the reader against directly comparing results between different sections. For convenience, all experimental configurations are tabulated in Appendix \ref{app:training-details}.


